\section{Setup and the exact law}\label{sec:law}

Let $q$ be a prime power, and let $V$ be a $2m$-dimensional vector space
over $\F_q$ with a nondegenerate alternating form of Gram matrix $J$.
An alternating matrix has zero diagonal and satisfies $K^{\tr}=-K$;
in characteristic two the zero-diagonal condition is essential.
Fix injective ordered frames
\[
 X_0:\F_q^a\longrightarrow V,\qquad B:\F_q^b\longrightarrow V,
\]
and put
\[
 H=X_0^{\tr}JX_0,\quad S=B^{\tr}JB,\quad
 \rank H=2h,\quad\rank S=2s.
\]
For $g$ uniform on $\Sp(V)$, define the compression
\begin{equation}\label{eq:compression}
 M=B^{\tr}JgX_0\in\Mat_{b,a}(\F_q).
\end{equation}
Using a covector frame instead is equivalent, since $J$ identifies $V$
with its dual. No assumption is made about the relative position of the
two frames.

For an alternating space of dimension $d$ and rank $2r$, define
\begin{equation}\label{eq:E}
 \mathcal E_m(d,r)=
 \prod_{i=0}^{r-1}\bigl((q^{2(m-i)}-1)q^{2(m-i)-1}\bigr)
 \prod_{i=0}^{d-2r-1}(q^{2(m-r)-i}-q^i)
\end{equation}
when $0\le2r\le d$ and $d-r\le m$, and set it to zero otherwise.
Empty products are one. This counts injective form-preserving maps into
$V$. On the feasible range it also equals
\begin{equation}\label{eq:Ecompact}
 \mathcal E_m(d,r)=
 q^{\,2mr-r^2+\binom{d-2r}{2}}
 \prod_{j=m-d+r+1}^{m}(q^{2j}-1).
\end{equation}
In particular $\mathcal E_m(0,0)=1$.

Let $\gb nk$ be the Gaussian coefficient, zero outside $0\le k\le n$,
and let
\[
 G_i(q)=|\GL_i(\F_q)|=\prod_{j=0}^{i-1}(q^i-q^j),\qquad G_0(q)=1.
\]
For $z,p,t\ge0$ with $p+t\le z$, put $w=z-p-t$ and set
\begin{equation}\label{eq:T}
 \mathcal T_{z,p,t}(k)=
 \sum_{\substack{0\le i\le\min(p,t,k)\\0\le k-i\le w}}
 \gb pi\gb ti G_i(q)\gb w{k-i}
 q^{(k-i)(p+t-i)}.
\end{equation}
This is a nonnegative integer, symmetric in $p,t$.

\Needspace{17\baselineskip}
\begin{theorem}[Exact arbitrary compression law]\label{thm:law}
For $C\in\Mat_{b,a}(\F_q)$, form
\begin{equation}\label{eq:joint}
 \mathcal G_C=\begin{pmatrix}S&C\\-C^{\tr}&H\end{pmatrix},
 \quad r=\tfrac12\rank\mathcal G_C,\quad z=a+b-2r.
\end{equation}
Put
\begin{equation}\label{eq:pt}
 p=\dim(\rad S\cap\ker C^{\tr}),\qquad
 t=\dim(\rad H\cap\ker C).
\end{equation}
Then, over every finite field,
\begin{equation}\label{eq:law}
 \boxed{\displaystyle
 \PP(M=C)=
 \frac{\sum_{k=0}^{z}\mathcal T_{z,p,t}(k)\,
                   \mathcal E_m(a+b-k,r)}
      {\mathcal E_m(b,s)\mathcal E_m(a,h)}.}
\end{equation}
The likelihood relative to uniform matrices is $q^{ab}\PP(M=C)$.
For fixed input dimensions and forms, each probability depends on $C$
only through $(r,p,t)$.
\end{theorem}

\begin{theorem}[Support and least ambient dimension]\label{thm:support}
In this setting,
\begin{equation}\label{eq:support}
 \boxed{\PP(M=C)>0\quad\Longleftrightarrow\quad r+\max(p,t)\le m.}
\end{equation}
More generally, for arbitrary alternating $S,H$ and arbitrary $C$, the
least dimension of a nondegenerate alternating space admitting
individually injective frames with these three Gram blocks is
$2\bigl(r+\max(p,t)\bigr)$.
\end{theorem}

The combined frame need not be injective: the source and target images
may intersect. Keeping only $k=0$ in \eqref{eq:law} would omit those
intersections and give incorrect probabilities and support.

\section{Proof of the exact law}\label{sec:proof}

\begin{lemma}[Embedding count and transitivity]\label{lem:embedding}
An alternating space of dimension $d$ and rank $2r$ embeds injectively
and isometrically into $V$ if and only if $d-r\le m$. Its number of
embeddings is $\mathcal E_m(d,r)$, and $\Sp(V)$ acts transitively on them.
\end{lemma}

\begin{proof}
Split the domain into $r$ symplectic pairs and a radical of dimension
$d-2r$. Once the nondegenerate part is embedded, the radical must map
to an independent isotropic frame in its $2(m-r)$-dimensional
orthogonal complement. Existence is equivalent to
$d-2r\le m-r$, which is the asserted condition.

In dimension $2N$, the first member of a symplectic pair has $q^{2N}-1$
choices, and a partner with pairing one has $q^{2N-1}$ choices.
Pass to their orthogonal complement and repeat. This gives the first
product in \eqref{eq:E}. After $i$ radical vectors are chosen, their
common perpendicular space has dimension $2(m-r)-i$. Exclude their
$q^i$-element span to keep the next vector independent. This gives the
second product.

An isometry between two embedded subspaces extends to the ambient
symplectic space by Witt's lemma \cite[Theorem 3.4]{SprehnWahl}. Directly,
split off the nondegenerate part, complete each radical basis to
hyperbolic pairs in its orthogonal complement, extend to symplectic
bases, and map one completed basis to the other. Thus transitivity
also covers degenerate subspaces and characteristic two.
\end{proof}

Consequently $gX_0$ is uniform on the $\mathcal E_m(a,h)$ injective
frames with Gram matrix $H$.

\begin{lemma}[Avoiding two disjoint subspaces]\label{lem:transverse}
Let $R$ have dimension $z$ and let $P,T\subset R$ be disjoint subspaces
of dimensions $p,t$. The number of $k$-subspaces $K\subset R$ with
$K\cap P=K\cap T=0$ is $\mathcal T_{z,p,t}(k)$. Such subspaces exist
exactly for
\begin{equation}\label{eq:kmax}
 0\le k\le z-\max(p,t).
\end{equation}
\end{lemma}

\begin{proof}
Fix $R=P\oplus T\oplus W$, with $\dim W=w$.
The subspace $K_0=K\cap(P\oplus T)$, of dimension $i$, must be the
graph of an isomorphism between an $i$-subspace of $P$ and an
$i$-subspace of $T$. There are $\gb pi\gb tiG_i(q)$ choices.
The projection of $K$ onto $W$ is a $(k-i)$-subspace $L$, with
$\gb w{k-i}$ choices. For fixed $K_0,L$, the remaining data are an
arbitrary map $L\to(P\oplus T)/K_0$, giving
$q^{(k-i)(p+t-i)}$ lifts. Such a lifted graph avoids both sides because
its vectors with zero $W$ projection belong to $K_0$. Summation proves
\eqref{eq:T}.

Injectivity of $K\to R/P$ and $K\to R/T$ bounds its dimension by
$z-\max(p,t)$. To attain this bound, take the graph of an isomorphism
between subspaces of dimension $\min(p,t)$ in the two sides and add
$W$. All smaller dimensions occur inside this maximum subspace.
\end{proof}

\begin{proof}[Proof of Theorem \ref{thm:law}]
Equip $D=\F_q^b\oplus\F_q^a$ with the form $\mathcal G_C$.
If a source frame $X$ realizes $C$, the combined map
$F=[B,X]:D\to V$ has this Gram matrix. Its kernel $K$ lies in
$R=\rad\mathcal G_C$. Injectivity of the individual frames is
equivalent to $K$ avoiding both coordinate summands.
Their intersections with $R$ are precisely the spaces of dimensions
$p,t$ in \eqref{eq:pt}. They are disjoint, so Lemma
\ref{lem:transverse} counts all possible kernels of each dimension.

For a fixed such $K$ of dimension $k$, the quotient $D/K$ has
dimension $a+b-k$ and rank $2r$. The form descends because $K\subset R$.
The target coordinate summand still embeds in this quotient.
An injective isometric embedding of $D/K$ into $V$, restricting to the
fixed target frame $B$, gives exactly one source frame with combined
kernel exactly $K$, and every such source frame arises this way.

If the quotient embeds, its embeddings restrict onto every target
frame with Gram matrix $S$: begin with one embedding and use
transitivity on target frames. The same group action identifies all
restriction fibers. Thus the number restricting to the fixed $B$ is
\[
 \frac{\mathcal E_m(a+b-k,r)}{\mathcal E_m(b,s)}.
\]
If the quotient does not embed the numerator is zero.
Sum over $K$ and divide by $\mathcal E_m(a,h)$ to obtain
\eqref{eq:law}. No separate basis choice for the intersection is needed.
\end{proof}

\begin{proof}[Proof of Theorem \ref{thm:support}]
The quotient for a $k$-dimensional kernel embeds if and only if
$a+b-k-r\le m$. The largest possible kernel dimension is
$z-\max(p,t)$. Some $k$ works exactly when
\[
 \max(0,a+b-r-m)\le z-\max(p,t).
\]
Since $z=a+b-2r$ and $z\ge p+t$, this is equivalent to
$r+\max(p,t)\le m$.
For the ambient-minimum assertion choose a maximum-dimensional
avoiding $K$. The quotient then has rank $2r$ and radical dimension
$\max(p,t)$, so it embeds in dimension $2(r+\max(p,t))$.
Both coordinate frames remain injective, and necessity excludes
every smaller dimension.
\end{proof}

\section{Consequences and special cases}\label{sec:consequences}

\subsection{Efficient exact evaluation}

The invariants are ordinary finite-field ranks:
\[
 r=\tfrac12\rank\mathcal G_C,\qquad
 p=b-\rank\begin{pmatrix}S\\C^{\tr}\end{pmatrix},\qquad
 t=a-\rank\begin{pmatrix}H\\C\end{pmatrix}.
\]
For $d=a+b$ they require $O(d^3)$ field operations.
The Gaussian coefficients and positive sum are evaluable with
polynomially many integer operations. Thus a matrix probability
requires no enumeration of the group or of the other matrices.
This is an arithmetic-operation statement; large-integer bit lengths
still matter.
The triple $(r,p,t)$ suffices for the likelihood. It need not classify
every stabilizer orbit of matrices, and the proof makes no such claim.

\subsection{Two isotropic frames}

For $S=H=0$ and $\rank C=v$, one has
$r=v$, $p=b-v$, $t=a-v$, and $z=p+t$. Then
\[
 \mathcal T_{z,p,t}(k)=\gb pk\gb tkG_k(q).
\]
The law depends only on ordinary matrix rank. Its support criterion
is $\max(a,b)\le m$, exactly individual feasibility. This explains
why rank is sufficient in the isotropic families of the \Article{}.

\subsection{Agreement with the Lagrangian-target formula}

Take $a=b=m$ and a Lagrangian target. A source frame in suitable
coordinates is $(A;B')$, with Gram equation
$A^{\tr}B'-B'^{\tr}A=H$.
For $K=\ker A$, $k=\dim K$, and $t=\dim(K\cap\rad H)$, the
\Article{} gives
\begin{equation}\label{eq:oldlag}
 \frac{\PP(A=C)}{q^{-m^2}}=
 \begin{cases}
 \displaystyle
 \frac{q^{\binom k2}\prod_{j=k-t+1}^{k}(1-q^{-j})}
 {\prod_{j=h+1}^{m}(1-q^{-2j})},&H|_{K\times K}=0,\\[6pt]
 0,&H|_{K\times K}\ne0.
 \end{cases}
\end{equation}
To see this directly, reduce $A$ to
$\operatorname{diag}(I_{m-k},0_k)$. The Gram equation fixes the cross
block of $B'$, the upper square block has
$q^{(m-k)(m-k+1)/2}$ choices, and the remaining block must be
injective on a $t$-space. This and Theorem \ref{thm:law} count the
same frames, so the formulas agree. The verifier also checks this
reduction in exact arithmetic.

\subsection{Beyond ordinary rank}

Take a nondegenerate plane as both source and target in dimension
four. In projection coordinates, likelihood distinguishes whether
the projected pair has pairing one. For $q\ge3$, invertible
$2$ by $2$ matrices can have either pairing one or another nonzero
pairing, so equal ordinary rank does not imply equal likelihood.
The next sections quantify this failure in unbounded dimensions.
